%----------------------------------
\section{Process Documentation}
\label{sec:Process Documentation}
%----------------------------------

This appendix provides evidence of the process conducted to produce this assignment: how the
Part 1 position paper was selected and evidenced, and how the Part 2 re-analysis of Task 1 was
carried out, including the code run, its raw console output, and a reflection on issues found
along the way.

\subsection{Part 1: sourcing and verifying the position paper and case study}

The five candidate position papers were checked against their publisher pages and, where
available, secondary reporting (e.g.\ The Conversation, university media coverage) to confirm
scope and findings before one was selected for review. \emph{Relational Futures} was chosen
because it is a substantive, evidenced report (36-respondent survey plus 22 yarning-circle
participants) with a clear purpose and audience, rather than a narrower methods paper. The case
study (predictive risk modelling in Australian child protection) was verified against a
peer-reviewed scoping review~\cite{krakouer2021analysing} rather than a secondary summary, and
was deliberately chosen to be outside the AI systems (companions, chatbots, generative tools)
that \emph{Relational Futures} itself discusses, per the assignment's requirement that the case
study not duplicate the paper's own examples.

\subsection{Part 2: reconstructing the Task 1 pipeline}

Task 1: Part 1 was written from an analysis whose code was not submitted with the report (only the
LaTeX write-up and result figures were retained). To evaluate that analysis's cross-validation
strategy, learning curves, and bias, the pipeline was rebuilt from the two original raw datasets
(UCI Online Retail II~\cite{chen2019onlineretailii} and UCI Online Shoppers Purchasing Intention
Dataset~\cite{sakar2018shoppers}) by following the preprocessing description in Task 1's own
process documentation (cleaning rules, product-week panel construction, feature list, split
dates). The reconstruction was verified against Task 1's reported figures before being extended:
it reproduced 106 weeks, a 31,800-row panel, 29,400 rows after lag construction, and 97.2\% of raw
rows retained after cleaning, all matching Task 1 exactly, and produced held-out metrics within a
few percentage points of Task 1's reported MAE/RMSE/WAPE and PR-AUC/ROC-AUC/F1 (Listing~\ref{lst:pipeline_a}
and Listings~\ref{lst:train_a}--\ref{lst:train_b}). This gave confidence that the reconstruction
was faithful enough to build the Part 2 analysis on.

\lstinputlisting[language=Python, basicstyle=\ttfamily\scriptsize, caption={Dataset A preprocessing: cleaning, product-week panel construction,
and feature engineering.}, label={lst:pipeline_a}]{code/pipeline_a.py}

\lstinputlisting[language=Python, basicstyle=\ttfamily\scriptsize, caption={Dataset A: temporal train/test split and model training (Random
Forest, Gradient Boosting, naive baseline).}, label={lst:train_a}]{code/train_a.py}

\textbf{Console output (Dataset A reconstruction):}
\begin{lstlisting}
raw (1067371, 8)
cleaned (1037007, 8) retained 97.155%
n weeks 106
panel shape (31800, 5)
after lag construction (29400, 16)
train (23400, 16) test (6000, 16)
Naive  MAE 127.07 RMSE 1076.38 WAPE 0.693
RF     MAE 112.03 RMSE 1068.95 WAPE 0.611
GB     MAE 112.45 RMSE 1068.85 WAPE 0.614
\end{lstlisting}

\lstinputlisting[language=Python, basicstyle=\ttfamily\scriptsize, caption={Dataset B: preprocessing, stratified train/test split, and model
training with F1-tuned threshold.}, label={lst:train_b}]{code/train_b.py}

\textbf{Console output (Dataset B reconstruction):}
\begin{lstlisting}
(12330, 18)
train (9247, 28) test (3083, 28) pos rate train 0.1548 test 0.1547
RF         PR-AUC 0.725 ROC-AUC 0.920 F1 0.661 (thr=0.372)
GB         PR-AUC 0.712 ROC-AUC 0.924 F1 0.659 (thr=0.338)
Always-no  PR-AUC 0.155 ROC-AUC 0.500 F1 0.000
\end{lstlisting}

\subsection{Cross-validation comparison}

\lstinputlisting[language=Python, basicstyle=\ttfamily\scriptsize, caption={Comparing naive random K-fold against TimeSeriesSplit (Dataset A) and
StratifiedKFold (Dataset B).}, label={lst:cv}]{code/cv_analysis.py}

\textbf{Console output:}
\begin{lstlisting}
=== Dataset A: naive random 5-fold CV (leaks future into past) ===
Random KFold WAPE per fold: [0.69, 0.749, 0.68, 0.694, 0.66] mean 0.695

=== Dataset A: TimeSeriesSplit 5-fold CV (respects chronology) ===
TimeSeriesSplit WAPE per fold: [0.572, 0.636, 0.716, 0.662, 0.605] mean 0.638

=== Dataset B: plain KFold (ignores 15.5% class imbalance) ===
Validation-fold positive rate (%): [16.7, 14.2, 15.2, 15.8, 15.5]
PR-AUC per fold: [0.758, 0.72, 0.754, 0.762, 0.761] mean 0.751 std 0.016

=== Dataset B: StratifiedKFold (preserves 15.5% class balance each fold) ===
Validation-fold positive rate (%): [15.5, 15.5, 15.5, 15.5, 15.5]
PR-AUC per fold: [0.761, 0.724, 0.757, 0.758, 0.735] mean 0.747 std 0.015
\end{lstlisting}

\subsection{Learning curve generation}

\lstinputlisting[language=Python, basicstyle=\ttfamily\scriptsize, caption={Learning curve generation for both datasets.}, label={lst:lc}]{code/learning_curves.py}

\textbf{Console output:}
\begin{lstlisting}
Dataset A learning curve (train size -> WAPE):
  n=  2340  WAPE=0.678
  n=  4680  WAPE=0.648
  n=  7020  WAPE=0.653
  n=  9360  WAPE=0.607
  n= 11700  WAPE=0.610
  n= 14040  WAPE=0.623
  n= 16380  WAPE=0.616
  n= 18720  WAPE=0.616
  n= 21060  WAPE=0.622
  n= 23400  WAPE=0.618

Dataset B learning curve (train size -> PR-AUC):
  n=   924  PR-AUC=0.641
  n=  1849  PR-AUC=0.690
  n=  2774  PR-AUC=0.705
  n=  3698  PR-AUC=0.712
  n=  4623  PR-AUC=0.715
  n=  5548  PR-AUC=0.721
  n=  6472  PR-AUC=0.710
  n=  7397  PR-AUC=0.719
  n=  8322  PR-AUC=0.721
  n=  9247  PR-AUC=0.720
\end{lstlisting}

\subsection{Bias analysis (Fairlearn MetricFrame equivalent)}

Fairlearn's package could not be installed in the sandboxed, offline analysis environment used to
produce this deliberation (no outbound package index access). Rather than substitute an unrelated
tool, its \texttt{MetricFrame} methodology --- computing a metric separately for each group defined
by a sensitive/grouping feature, and comparing the spread --- was implemented directly with
\texttt{pandas.groupby} and scikit-learn metric functions, as shown below.

\lstinputlisting[language=Python, basicstyle=\ttfamily\scriptsize, caption={Disaggregated (Fairlearn MetricFrame-equivalent) bias analysis for
both datasets.}, label={lst:bias}]{code/bias_analysis.py}

\textbf{Console output:}
\begin{lstlisting}
=== Dataset A: disaggregated error by product revenue tier ===
                     n_obs         MAE      WAPE  mean_actual_demand
tier
Low (bottom third)  2000.0  133.297629  0.713058            186.9380
Mid (middle third)  2000.0   82.728300  0.631494            131.0040
High (top third)    2000.0  121.321784  0.523479            231.7605

=== Dataset B: disaggregated error by VisitorType ===
                   n_sessions  actual_conv_rate       FNR       FPR
VisitorType
New_Visitor             429.0          0.256410  0.245455  0.050157
Other                    13.0          0.230769  0.000000  0.000000
Returning_Visitor      2641.0          0.137827  0.354396  0.072025

=== Dataset B: disaggregated error by Region ===
        n_sessions  actual_conv_rate  selection_rate       FNR       FPR
Region
1           1258.0          0.162957        0.163752  0.321951  0.063628
2            269.0          0.171004        0.200743  0.304348  0.098655
3            566.0          0.146643        0.160777  0.373494  0.080745
4            303.0          0.158416        0.188119  0.250000  0.082353
5             70.0          0.142857        0.200000  0.200000  0.100000
6            204.0          0.132353        0.102941  0.518519  0.045198
7            188.0          0.154255        0.132979  0.310345  0.031447
8            110.0          0.081818        0.109091  0.444444  0.069307
9            115.0          0.173913        0.182609  0.200000  0.052632
\end{lstlisting}

\subsection{Reflection}

The main process risk in Part 2 was that the original Task 1 code was unavailable, so the
reconstruction could have quietly drifted from what was actually run in Task 1 without the author
noticing. This was mitigated by validating every intermediate quantity (rows retained, panel
shape, post-lag row count, held-out metrics) against Task 1's own reported numbers before treating
the reconstruction as a basis for new analysis, and by explicitly disclosing the one place a new
issue was found (the \texttt{price\_rel} look-ahead feature) rather than silently patching it. The
bias analysis substituted a manually implemented equivalent for Fairlearn because the package
could not be installed offline; this substitution is disclosed in Section~4.3 of the main text
rather than presented as if Fairlearn itself had been used.

%----------------------------------
\section{Link to GitHub Repository}
%----------------------------------

\url{https://github.com/Sahanaa-N/csds.git}
